% BHU PYQ -- Professional Exam Template
\documentclass[11pt,a4paper]{article}

\usepackage[a4paper,top=2.5cm,bottom=2.5cm,left=2.8cm,right=2.8cm,headheight=14pt]{geometry}
\usepackage{amsmath,amssymb}
\usepackage[T1]{fontenc}
\usepackage[utf8]{inputenc}
\usepackage{lmodern}
\usepackage{microtype}
\usepackage{fancyhdr}
\usepackage{enumitem}
\usepackage{listings}
\usepackage{hyperref}
\usepackage{lastpage}
\usepackage{parskip}

\hypersetup{colorlinks=true,urlcolor=black,linkcolor=black,
  pdftitle={BVCA-404: Microprocessor -- BHU PYQ}}

\pagestyle{fancy}
\fancyhf{}
\renewcommand{\headrulewidth}{0.4pt}
\renewcommand{\footrulewidth}{0.4pt}
\fancyhead[L]{\small   \textit{Banaras Hindu University}}
\fancyhead[R]{\small   \textit{BVCA-404: Microprocessor}}
\fancyfoot[C]{\small Page  \thepage\ of \pageref{LastPage}}

\lstset{
  basicstyle=\small tfamily,
  frame=single,
  breaklines=true,
  columns=flexible,
  keepspaces=true
}


\newlist{parts}{enumerate}{2}
\setlist[parts,1]{label=(\alph*),leftmargin=*,itemsep=4pt,topsep=3pt}
\setlist[parts,2]{label=(\roman*),leftmargin=*,itemsep=2pt,topsep=2pt}

\newcommand{\pts}[1]{\hfill   \textit{[#1]}}


\begin{document}

\begin{center}
  {\large\bfseries BANARAS HINDU UNIVERSITY}\\[2pt]
  {
\normalsize B.Voc. Semester IV Examination, 2022-23}\\[6pt]
  \rule{0.8\linewidth}{0.5pt}\\[6pt]
  {\large\bfseries Paper: BVCA-404 --- Microprocessor}\\[4pt]
  {
\normalsize Time: Three Hours \hfill Maximum Marks: 70}
\end{center}

\rule{\linewidth}{0.4pt}

\smallskip



\noindent   \textit{Instructions: Answer the questions in each section as instructed. All questions carry marks as indicated.}

\bigskip




\noindent {\large\bfseries SECTION --- A}
\rule{\linewidth}{0.4pt}\smallskip

Long-Answer Questions. Answer any two questions: \hfill   \textbf{2 $ \times$ 12 = 24 Marks}

\medskip



\noindent   \textbf{Question 1.} Explain the architecture of the 8085 microprocessor, describing each component and how they interact with each other.

\medskip



\noindent   \textbf{Question 2.} Discuss the various types of addressing modes in the 8085 microprocessor, providing a suitable example for each.

\medskip



\noindent   \textbf{Question 3.} Explain the role of the stack and the stack pointer in the 8085 microprocessor. Illustrate your explanation with practical examples.

\medskip



\noindent   \textbf{Question 4.} Describe the differences between a microprocessor and a microcontroller, with emphasis on the 8085 as a specific example of a microprocessor.

\bigskip


\noindent {\large\bfseries SECTION --- B}
\rule{\linewidth}{0.4pt}\smallskip




\noindent   \textbf{Question 5.} Attempt any six parts (Answer in about 200 words each): \hfill   \textbf{6 $ \times$ 6 = 36 Marks}

\begin{parts}
  \item Briefly explain the function of the Arithmetic Logic Unit (ALU) in the 8085 microprocessor.
  \item List and explain the status of any three flags in the 8085 microprocessor's flag register.
  \item Define the terms ``Opcode'' and ``Operand'' with respect to the 8085 microprocessor's instruction format.
  \item Explain the purpose and classifications of the instruction set in a microprocessor.
  \item If a program starts at memory location   \texttt{2000H}, and the first instruction is   \texttt{LDA 2050H}, what will be the content of the accumulator and the program counter after this instruction is executed?
  \item Write an assembly language program for the 8085 microprocessor to add two 8-bit numbers stored at memory locations   \texttt{3000H} and   \texttt{3001H}, and store the result in memory location   \texttt{3002H}.
  \item Given that the stack pointer (SP) is initially at   \texttt{4000H}, what will be its value after executing the following sequence of instructions?
  
\begin{lstlisting}[language={[x86masm]Assembler}]
PUSH H
PUSH D
POP H
  \end{lstlisting}
  \item Explain the memory address accessed by the instruction sequence:
  
\begin{lstlisting}[language={[x86masm]Assembler}]
LXI H, 4050H
MOV A, M
  \end{lstlisting}
  \item What is interrupt-driven I/O and how is it implemented in the 8085 microprocessor?
\end{parts}

\bigskip


\noindent {\large\bfseries SECTION --- C}
\rule{\linewidth}{0.4pt}\smallskip




\noindent   \textbf{Question 6.} Attempt all parts: \hfill   \textbf{10 $ \times$ 1 = 10 Marks}

\begin{parts}
  \item What does the acronym ALU stand for?
  \item Name the register that stores the result of an arithmetic or logical operation in the 8085 microprocessor.
  \item What is the size of the data bus in the 8085 microprocessor?
  \item Which pin on the 8085 microprocessor is used for a non-maskable hardware interrupt?
  \item How many machine cycles does the register-to-register   \texttt{MOV} instruction require?
  \item How do you initialize the stack pointer in the 8085 microprocessor?
  \item What does the   \texttt{HLT} instruction do?
  \item Which flag is set if an arithmetic or logical operation results in zero in the 8085 microprocessor?
  \item What is the opcode (in hex) for the   \texttt{NOP} instruction in the 8085 microprocessor?
  \item How many bytes does the   \texttt{LDA} instruction occupy?
\end{parts}

\vfill
\rule{\linewidth}{0.4pt}

\begin{center}\small
\href{https://bhu-exam.vercel.app/}{Click here to visit our website}\\[4pt]\href{https://me-aryan.vercel.app/}{Click here to visit our contributor}\\[4pt]\href{https://quashan-me.vercel.app/}{Click here to visit our contributor}
\end{center}
\end{document}
